\subsection{Where the Democratic Dividend Actually Comes From}
\label{sec:10.1.1}
AI that flags disinformation, exposes corruption, protects an election's integrity, or helps a citizen engage with government: these are real applications, and none of them is what makes AI protect democratic values rather than undermine them. Section 8.7.6 names the thing that does: what it calls a democratic dividend, available to a state that already has functioning rule-of-law institutions, a free press, and courts a citizen can actually use, and unavailable to a state that lacks those, regardless of which application gets built on top of identical technology.

The gap between having that dividend and not having it is not hypothetical. Freedom House's 2023 assessment of internet freedom found at least twenty-two countries where law already requires or incentivizes a platform to deploy machine learning against content the state disfavors, almost never described in those terms — the operative language is national security or a threat of disinformation, not censorship \autocite{funk2023freedom}. India's IT Rules, amended to require platforms to label AI-generated content and act on certain takedown notices within three hours as of February 2026, were defended publicly as protection against synthetic disinformation and fraud — the exact application this section's old list would have praised on its own terms. A lawyer who has litigated India's free-speech cases has documented that enforcement in practice concentrates on parody and criticism of the prime minister, not the harms the rules name \autocite{gupta2026three}. The application did not fail at what it was built to do. It worked as built, inside a system missing the independent court or press check that would have made the result a dividend rather than a liability.

None of the applications on the old list is wrong to want. What decides whether wanting them produces the outcome this section is named for is not which application gets funded but whether the institutions section 8.7.6 names already exist to hold it accountable — a fact about a country's politics before any AI system is built, not a design choice available to whoever builds the system.
